% tab.tex: for Software: Practice and Experience (Wiley NJD class;
% WileyNJD-v2.cls, NJDnatbib.sty and WileyNJD-AMA.bst sit in this
% directory -- they are not on CTAN, so tectonic cannot fetch them;
% everything else tectonic pulls from its bundle on demand).
% Code arrives via snips/, cut fresh from ../tab's sources by tab's
% ,/snips (see Makefile): nothing here is pasted, nothing can rot.
% Bib entries live at the end of this file (filecontents), so this
% paper's whole footprint is this file, the Makefile and the three
% Wiley files above.
\documentclass[AMA,STIX2COL]{WileyNJD-v2}
\usepackage{listings,xcolor}
\articletype{Experience Report}

% Code style after the ezr paper (tosem10.tex): very small type, a
% hairline left rule, line numbers every 5th -- and numbering that
% runs CONTINUOUSLY across listings, so the paper carries one set of
% line numbers for the whole toolkit.
\newcommand{\vtiny}{\fontsize{8}{9}\selectfont}
\lstset{language={},
  basicstyle=\vtiny\ttfamily, breaklines=true,
  postbreak=\mbox{\textcolor{gray}{$\hookrightarrow$}\,},
  commentstyle=\color{gray}\itshape,
  keywordstyle=\bfseries,
  morekeywords={function,local,return,end,if,then,else,elseif,
                for,while,do,in,and,or,not,nil,true,false},
  morecomment=[l]{--},
  columns=fullflexible, keepspaces=true,
  numbers=left, firstnumber=last, stepnumber=5,
  numberstyle=\fontsize{4}{4}\selectfont\color{black!50},
  numbersep=5pt, xleftmargin=8pt,
  frame=leftline, framerule=0.3pt, rulecolor=\color{black!25},
  aboveskip=1ex, belowskip=1ex}
% \snip{NUM.add} pulls the code AND makes Listing~\ref{lst:NUM.add}
% (and its line numbers) citeable from the prose.
\newcommand{\snip}[1]{\lstinputlisting[caption={\texttt{#1}},
  label={lst:#1}]{snips/#1.lua}}

\begin{document}
\title{Use the model, ship the program:\\
  neuro-guided synthesis of very fast learners}
\author[1]{Tim Menzies}
\authormark{MENZIES}
\address[1]{\orgdiv{Computer Science}, \orgname{North Carolina State
  University}, \orgaddress{\state{NC}, \country{USA}}}
\corres{Tim Menzies. \email{timm@ieee.org}}

\abstract[Summary]{On [N] tabular SE data sets, a Naive Bayes
classifier over seven discretized bins comes within [X]\% of LightGBM
and TabICL (by recall and false alarm).  The whole learner is [N]
lines of Lua with no dependencies; it trains and tests in [T] seconds
on a laptop.  We describe its design, the experience of keeping it
small, and the working style it enables: many fast runs, and only an
occasional slow check against an expensive reference.}

\keywords{naive Bayes, discretization, baselines, Lua, little
languages, reproducibility}

\maketitle

\section{Introduction}

MSR May 19, 2013

Every month brings a bigger model.  Every month, last month's results
look old.

This paper offers a result that will not age so fast.  On [N] tabular
SE data sets, a Naive Bayes classifier with seven discretized bins
comes within [X]\% of gradient boosting (GBM, LightGBM) and TabICL,
measured by recall (pd) and false alarms (pf).  The whole learner is
[N] lines of Lua.  It runs in [T] seconds on a laptop.  A student can
read it in an afternoon.

After two years of intensive work with LLMs, writing papers on
software analytics, we assert that science-oriented programming is
different from commercial programming.  Commercial code optimizes
for delivery: ship the feature, minimize time and cost, and let
size and opacity hide behind test suites and teams.  Scientific
code optimizes for something else -- the veracity of a conclusion.
Here, the program is not a product that supports the argument; the
program \emph{is} the argument, and every line is a claim that a
stranger must be able to check.

That difference makes minimality an end in itself, not an
aesthetic.  A scientific deliverable is judged by how quickly a
skeptical reader reaches either ``I believe this'' or ``here is
your bug.''  So we practice \emph{seeking minima as a
methodology}: delete until nothing can hide.  Seeking minima is
not simple, since a minimum means little until it is benchmarked
against the maximal alternative -- hence every small learner here
is scored against LightGBM and TabICL, so the cost and the benefit
of smallness are both on the table.  The reward is what the minima
teach: each dramatic simplification found below (one update
primitive, training inside constructors, learners as two or three
questions asked of shared summaries) was a thing we did not know
until the surrounding code was gone.

There is also a lesson here about building neuro-symbolic systems.
This kit was written with an LLM used as a sounding board, a tool
for exploring options, a research assistant doing the mechanical
work -- while a small symbolic artifact (a few hundred lines of
Lua, pinned by golden tests) disciplined every suggestion: any
change had to reproduce byte-identical output and hold the
baseline speed, or it died.  Many fast probes, aggressive
deletion, and an executable referee between the two.  We offer
that loop as one answer to how to do science in the age of LLMs.

We do not claim that big tools are useless.  We claim that most of
the time, they are not needed.  Science works best with many fast
experiments and only a few slow ones.  So we should run many cheap
probes and check against an expensive model only now and then.

This is an old idea.  In the 1960s, information systems buried
managers in reports.  Decision support systems replied with a
question: what is the least information that most changes a decision?
AI now buries us in code.  We ask that question again, and we answer
it in a small language.

We wrote this for three readers:
\begin{itemize}
\item Scientists who want results they can repeat on a laptop, and
  that will still hold after the next model release.
\item Teachers who need code students can read end to end.
  Schr\"odinger warned that if you cannot, in the long run, tell
  everyone what you have been doing, ``your doing has been
  worthless''~\cite{schrodinger1951science}.  Small code makes that
  telling possible.
\item Lua users who already believe a little language, used well,
  goes a long way~\cite{ierusalimschy1996lua}.  Here is evidence.
\end{itemize}

Our contributions are:
\begin{itemize}
\item An empirical result: a 7-bin Naive Bayes is close to
  state-of-the-art on pd/pf for [N] data sets.
\item A method: many fast runs, one slow check.
\item A tool: tab.lua, [N] lines, no dependencies.
\item An argument: some results outlast the current model cycle, and
  we should look for them.
\end{itemize}

\section{Motivation}

\subsection{Overload, again}

In 1938, Buckminster Fuller described ephemeralization: doing ``more
and more with less and less''~\cite{fuller1938nine}.  Technology
delivered that promise, then delivered too much.  Gross gave the
problem a name in 1964~\cite{gross1964managing}.  Toffler made it
famous in 1970~\cite{toffler1970future}.  In between, Ackoff saw it
inside companies.  Managers did not lack relevant information.  They
drowned in irrelevant information~\cite{ackoff1967misinformation}.
Simon put it in one line: ``a wealth of information creates a poverty
of attention''~\cite{simon1971designing}.

The answer then was decision support~\cite{gorry1971framework}.  Do
not give people every report.  Give them the smallest thing that
changes what they do next.

LLMs repeat the cycle.  Code and text now cost almost nothing to
produce.  Reading them, checking them, and trusting them still cost a
lot.  So the problem has moved from making things to choosing among
them, and the right response is the old one: find the minimum that
matters.

\subsection{Programming under examination}

Repetition matters because we all make mistakes.  In 2013 we had a
paper accepted at MSR on cross-company defect prediction.  It beat
the prior state of the art.  Within three weeks, before the
conference started, a graduate student trying to reproduce it found a
coding error: a bad re-insert of class labels after a distance
calculation.  Once we fixed it, the result went away.  We reported
this at the conference.  That error was found because the code was
small enough to read, and the data was open enough to rerun.  Large,
opaque, expensive pipelines hide errors like this.  Small, cheap ones
expose them.

That episode set a working rule we have kept since: offer
everything.  All the code, all the data, the seeds, the splits, the
scripts that drew every figure -- so that others can repeat,
improve, or refute the result at the cost of one download.  This is
partly service and partly discipline.  Write every line expecting
that it will be examined, in great detail, by some researcher
somewhere in the world, and the line comes out different: shorter,
plainer, with its assumptions on the surface.  Reputation in
science is not built on being right the first time; it is built on
being checkable, and on reporting quickly when the check fails.

Big models work against this.  A result that needs a GPU cluster
can be checked only by people who have one.  Most of the world's
scientists, students, and small labs do not.  For them, current AI
research is something they read about but cannot repeat.  When
results cannot be repeated by most people, it is hard to call them
science.

\subsection{What the logs say}

The difference can be measured.  This kit was built with an LLM
assistant, and the assistant's transcripts record every request
made of it.  Over two months and 868 prompts (this kit, plus four
sibling research projects), the requests sort as follows: 105
asked for renames or uniformity (``is \texttt{i} only ever used
for self?'', ``constructors before methods''); 72 interrogated
code that already worked (``why do we need \texttt{leaves}?'',
``surely only \texttt{Num} needs \texttt{w}?''); 61 demanded
verification (golden tests, byte-identical output, ``same job at
same speed as the lua we started with''); 57 asked for something
to be deleted, folded, or shrunk.  Not one asked for faster
delivery.  (The classifier is \texttt{etc/logmine.py} in this
paper's repository: thirty lines that read any Claude Code
transcript directory, so readers can audit our audit -- or their
own practice.)  A commercial reviewer would read most of that traffic
as churn.  Here it is the work: the majority of the editing budget
went to making the artifact smaller, more uniform, and more
interrogated, under a regime where behavior and speed were frozen
by tests.  Form is content when the artifact is the argument.

\subsection{The ground keeps moving}

There is a third problem.  Whatever we build on today's model, the
next model will replace.  A paper tuned to one release has a short
shelf life.  Researchers end up on a treadmill, spending effort to
keep up instead of to understand.

So we look for results that do not depend on the current release.
Some ideas have lasted decades: discretize, count, compare, keep the
simplest model that works.  A 7-bin Naive Bayes will mean the same
thing in ten years as it does today.  Its baseline does not expire.
When a new model arrives, the right question is still ``how much
better than the simple thing?''  And often the honest answer is ``not
much.''

\subsection{Little languages, fast runs}

Our tool choice follows from this.  Lua is a small language with a
small core and few dependencies~\cite{ierusalimschy1996lua}.  It
changes slowly, so code written years ago still runs.  It encourages
short scripts that do one thing.

That suits a research style with two speeds:
\begin{itemize}
\item Fast: thousands of cheap runs with small learners to explore
  the option space.
\item Slow: a few runs against an expensive reference (GBM, LightGBM,
  TabICL) to confirm what the fast runs found.
\end{itemize}

Most of the time goes into the fast loop.  The slow loop is a check,
not the main tool.

\section{Columns}

Columns summarize streams: counts for symbols, Welford's running mean
and standard deviation for numbers.

\snip{Sym}
\snip{SYM.add}
\snip{Num}
\snip{NUM.add}

\section{Naive Bayes}

Training is the constructor: one sub-table per klass.

\snip{Nb}
\snip{NB.like}
\snip{NB.guess}

\section{Range means}

Distance to heaven labels a row; bins collect those labels; a guess
is the variance-weighted, shrunken mean of the bins a row falls into.

\snip{TAB.ydist}
\snip{Rmeans}
\snip{RMEANS.guess}

\nocite{*}
\bibliography{refs}

% ---- bib entries, loaded at the end of this file -------------------
\begin{filecontents}[overwrite]{refs.bib}
@misc{menzies26tab,
  author = {Tim Menzies},
  title  = {tab: fast inference over tables},
  year   = {2026},
  url    = {https://github.com/timm/tab}}
@article{ierusalimschy1996lua,
  author  = {Ierusalimschy, Roberto and de Figueiredo, Luiz Henrique
             and Celes, Waldemar},
  title   = {Lua---an extensible extension language},
  journal = {Software: Practice and Experience},
  volume  = {26}, number = {6}, pages = {635--652}, year = {1996}}
@book{schrodinger1951science,
  author    = {Schr{\"o}dinger, Erwin},
  title     = {Science and Humanism},
  publisher = {Cambridge University Press}, year = {1951}}
@book{fuller1938nine,
  author    = {Fuller, R. Buckminster},
  title     = {Nine Chains to the Moon},
  publisher = {J. B. Lippincott}, year = {1938}}
@book{gross1964managing,
  author    = {Gross, Bertram M.},
  title     = {The Managing of Organizations},
  publisher = {Free Press of Glencoe}, year = {1964}}
@book{toffler1970future,
  author    = {Toffler, Alvin},
  title     = {Future Shock},
  publisher = {Random House}, year = {1970}}
@article{ackoff1967misinformation,
  author  = {Ackoff, Russell L.},
  title   = {Management misinformation systems},
  journal = {Management Science},
  volume  = {14}, number = {4}, pages = {B147--B156}, year = {1967}}
@incollection{simon1971designing,
  author    = {Simon, Herbert A.},
  title     = {Designing organizations for an information-rich world},
  booktitle = {Computers, Communications, and the Public Interest},
  publisher = {Johns Hopkins Press}, pages = {37--72}, year = {1971}}
@article{gorry1971framework,
  author  = {Gorry, G. Anthony and {Scott Morton}, Michael S.},
  title   = {A framework for management information systems},
  journal = {Sloan Management Review},
  volume  = {13}, number = {1}, pages = {55--70}, year = {1971}}
\end{filecontents}
\end{document}
